\section{Related Work}
\label{sec:related}

\subsection{Ultrasound-based liver fibrosis grading}

Automated fibrosis assessment from ultrasound has been pursued along several
lines. One line augments B-mode input with additional physical measurements,
such as contrast-enhanced micro-flow cines~\cite{microflow2023} or
elastography-derived stiffness~\cite{swe2018staging}, trading equipment
requirements for signal quality. A second line works from conventional B-mode
images alone, using generative augmentation to compensate for limited
data~\cite{usfibrosisgan2022}, higher-frequency acquisition to expose
finer texture~\cite{hfus2025}, or paired organ views to exploit relative
appearance~\cite{liverspleen2025}.

These methods almost universally target a small number of discrete stages of
fibrosis arising from viral, alcoholic or metabolic liver disease. The
schistosomal setting differs in both the target and the phenotype. The septal,
network-like fibrosis induced by \emph{S.~japonicum} presents differently from
the diffuse parenchymal change typical of other
etiologies~\cite{schistoelasto2018,schistopocus2024}, and the grading scale used
in endemic screening programs is finer than a three- or four-way stage decision.
SFibAI addressed this setting by casting grading as a posterior over 36 ordinal
levels and reporting its expectation~\cite{xu2026sfibai}; the present work
inherits that formulation and asks what anatomical structure adds to it.

\subsection{Ordinal and fine-grained prediction}

Treating graded severity as unordered classes discards the ordering that makes
grades meaningful, and a substantial literature addresses this. One family
reformulates the ordinal problem as a set of coupled binary decisions, enforcing
rank consistency either architecturally~\cite{coral2020} or through conditional
probability factorization~\cite{corn2023}. A second family discretizes a
continuous quantity into ordered bins and predicts a per-bin probability
distribution, as in ordinal depth
estimation~\cite{dorn2018}. A third family regularizes the predicted
distribution to be unimodal so that probability mass concentrates around
neighbouring levels rather than scattering across the
scale~\cite{unimodalbeta2021}.

In medical imaging, ordinal formulations have been applied to disease severity
grading, sometimes jointly with segmentation~\cite{sonnet2022}, and
uncertainty-aware variants have been developed for retinal
grading~\cite{drgraduate2020}. Our grading head follows the second family: a
36-level posterior whose expectation is the reported score, combined with a
neighborhood-aware soft-label term and a boundary-sensitive penalty inherited
from SFibAI~\cite{xu2026sfibai}. We do not claim a new ordinal formulation;
the fine-grained scale is the task setting within which the anatomical priors
are studied, and it matters here because subtle level distinctions are precisely
where additional structural context could plausibly contribute.

\subsection{Anatomy-informed and auxiliary-supervised learning}

The closest body of work injects anatomical knowledge into a network that
performs a different primary task. Several strategies recur.

\textbf{Attention conditioned on anatomy.} Attention-gated networks learn to
emphasize salient regions and suppress irrelevant background without explicit
localization supervision~\cite{attentiongated2019}. Anatomy-XNet extends this
by supervising attention with anatomical segmentation masks so that thoracic
disease classification attends to lung and heart
regions~\cite{anatomyxnet2022}. SynAP-Fib shares the intent of steering the
representation toward anatomically meaningful evidence, but its lesion signal is
a coarse box union rather than an organ segmentation, and the resulting
attention is injected into grading as a gated residual rather than used to
reweight the backbone's own spatial attention.

\textbf{Anatomy-partitioned architectures.} AGMB-Transformer routes features
through anatomy-defined branches for automated evaluation of root canal
therapy~\cite{agmbtransformer2021}, and anatomy-guided cascades exploit the
known containment of the optic cup within the optic disc~\cite{opticdisc2020}.
HoVer-Trans encodes the horizontal and vertical tissue layering of breast
ultrasound directly into the transformer's attention
structure~\cite{hovertrans2023}. These designs commit the architecture to a
fixed anatomical decomposition. SynAP-Fib instead keeps a single shared
backbone and treats anatomical position as a \emph{predicted} six-way posterior
that modulates grading, so no branch is tied to a particular position and no
position metadata is needed at inference.

\textbf{Multi-task learning with auxiliary heads.} Joint segmentation and
classification has been studied in breast ultrasound, where joint learning
improves outcomes of both
tasks~\cite{zhou2020mtlabus,shamtl2021}, and joint pose estimation with
tissue segmentation has been used for fetal
brains~\cite{petsnet2023}. Weakly supervised localization achieves a related
effect with far cheaper labels, deriving lesion evidence from image-level
supervision~\cite{covidweak2020,weakbreast3d2019}. In these designs the
auxiliary branch typically shares gradients with the primary task, so its
benefit and its interference are entangled: the auxiliary loss shapes the shared
representation directly. SynAP-Fib deliberately separates the two pathways with
a bidirectionally detached gradient boundary, so each branch contributes its
\emph{output} as context rather than its gradient as regularization.

\textbf{Positioning.} Individually, none of these ingredients is new:
anatomy-supervised attention, position-conditioned prediction, auxiliary
multi-task heads, weak-box supervision and gated feature fusion all have
precedent. Several of the closely related studies considered here use a
single auxiliary signal in their favourable configuration. The $2\times2$
comparison we report addresses a different question: whether anatomical
priors of distinct kinds compose, and it yields an answer that a
single-prior study could not have surfaced.
